% GENERATED FILE. Edit canonical lessons, then run npm run generate:books.
% canonical-source-hash: fnv1a64:48cc2de5d0e38869
% canonical-lessons: BN-C349-protibad-kora, BN-C349-oborodh, BN-C349-jonomot, BN-C349-niti, BN-C349-jonosongkhya

\chapter{Raising a Voice}
\label{ch:bn-raising-a-voice}
\hlchaptermodality{349}

\begin{chapteropening}
\textbf{By the end of this chapter:} \emph{I can say what people protest about and what they want.}

Complete the last lesson of chapter 349: I can say what people protest about and what they want.
\end{chapteropening}

\section[\texorpdfstring{\bn{প্রতিবাদ} \bn{করা} (prôtibād kôrā)}{প্রতিবাদ করা (prôtibād kôrā)}]{\bn{প্রতিবাদ} \bn{করা} (prôtibād kôrā) — to protest}

\label{lesson:BN-C349-protibad-kora}



\begin{quote}
Before the new one: say the Bengali for to collect dues, then the Bengali for to sanction.
\end{quote}

\subsection*{You'll want to know: \bn{প্রতিবাদ} \bn{করা}}
\textbf{\bn{প্রতিবাদ} \bn{করা}} — \emph{prôtibād kôrā} — \textquotedblleft{}to protest\textquotedblright{}.

\begin{quote}
\bn{ছাত্ররা} \bn{নতুন} \bn{নিয়মের} \bn{প্রতিবাদ} \bn{করল।}
\end{quote}

\begin{quote}
\emph{chhātrôrā nôtun niyômer prôtibād kôrlô.}
\end{quote}

\begin{quote}
The students protested against the new rule.
\end{quote}

Past (he/she …): \bn{তিনি} \bn{অন্যায়ের} \bn{প্রতিবাদ} \bn{করলেন}, \emph{tini ônyāyer prôtibād kôrlen}.

\begin{cousinweb}
\textbf{\bn{প্রতিবাদ} \bn{করা}} (\emph{prôtibād kôrā}) is Sanskrit \emph{prativāda}, speaking against, contradiction (Monier-Williams), with \bn{করা}. What is protested takes -\bn{এর}.
\end{cousinweb}

\begin{grammarlens}[title={What you've built}]
One more word for giving reasons and opinions.
\end{grammarlens}

\subsection*{Your turn}
\begin{itemize}
\raggedright
  \item \emph{Say it:} \emph{prôtibād kôrā}
  \item \emph{Say it:} \emph{prôtibād kôrā}, once more
  \item \emph{Say it:} \emph{chhātrôrā nôtun niyômer prôtibād kôrlô.} — The students protested against the new rule.
  \item \emph{Recall:} say the Bengali for to collect dues, then the Bengali for to sanction, then say \emph{prôtibād kôrā} again
\end{itemize}

\begin{tcolorbox}[breakable,colback=teal!4,colframe=teal!35!black,title={Before you move on}]
What does \bn{প্রতিবাদ} \bn{করা} mean? (\textquotedblleft{}To protest\textquotedblright{}.) Say it once more.
\end{tcolorbox}
\section[\texorpdfstring{\bn{অবরোধ} (ôbôrodh)}{অবরোধ (ôbôrodh)}]{\bn{অবরোধ} (ôbôrodh) — a blockade; a sit-in (blocking a road)}

\label{lesson:BN-C349-oborodh}



\begin{quote}
Before the new one: say the Bengali for to sanction, then the Bengali for to protest.
\end{quote}

\subsection*{You'll want to know: \bn{অবরোধ}}
\textbf{\bn{অবরোধ}} — \emph{ôbôrodh} — \textquotedblleft{}a blockade; a sit-in (blocking a road)\textquotedblright{}.

\begin{quote}
\bn{অবরোধের} \bn{জন্য} \bn{রাস্তায়} \bn{বাস} \bn{নেই।}
\end{quote}

\begin{quote}
\emph{ôbôrodher jônyô rāstāy bās nei.}
\end{quote}

\begin{quote}
There are no buses on the road because of the blockade.
\end{quote}

\begin{cousinweb}
\textbf{\bn{অবরোধ}} (\emph{ôbôrodh}) is Sanskrit \emph{avarodha}, hindrance or obstruction (Monier-Williams). In Bengali news it is a protest that blocks a road or a railway.
\end{cousinweb}

\begin{grammarlens}[title={What you've built}]
One more word for giving reasons and opinions.
\end{grammarlens}

\subsection*{Your turn}
\begin{itemize}
\raggedright
  \item \emph{Say it:} \emph{ôbôrodh}
  \item \emph{Say it:} \emph{ôbôrodh}, once more
  \item \emph{Say it:} \emph{ôbôrodher jônyô rāstāy bās nei.} — There are no buses on the road because of the blockade.
  \item \emph{Recall:} say the Bengali for to sanction, then the Bengali for to protest, then say \emph{ôbôrodh} again
\end{itemize}

\begin{tcolorbox}[breakable,colback=teal!4,colframe=teal!35!black,title={Before you move on}]
What does \bn{অবরোধ} mean? (\textquotedblleft{}A blockade; a sit-in (blocking a road)\textquotedblright{}.) Say it once more.
\end{tcolorbox}
\section[\texorpdfstring{\bn{জনমত} (jônômôt)}{জনমত (jônômôt)}]{\bn{জনমত} (jônômôt) — public opinion}

\label{lesson:BN-C349-jonomot}



\begin{quote}
Before the new one: say the Bengali for to protest, then the Bengali for a blockade.
\end{quote}

\subsection*{You'll want to know: \bn{জনমত}}
\textbf{\bn{জনমত}} — \emph{jônômôt} — \textquotedblleft{}public opinion\textquotedblright{}.

\begin{quote}
\bn{জনমত} \bn{সরকারের} \bn{বিরুদ্ধে।}
\end{quote}

\begin{quote}
\emph{jônômôt sôrkārer biruddhe.}
\end{quote}

\begin{quote}
Public opinion is against the government.
\end{quote}

\begin{cousinweb}
\textbf{\bn{জনমত}} (\emph{jônômôt}) is a modern compound of Sanskrit \emph{jana}, people, and \emph{mata}, an opinion (Monier-Williams): what the public thinks.
\end{cousinweb}

\begin{grammarlens}[title={What you've built}]
One more word for giving reasons and opinions.
\end{grammarlens}

\subsection*{Your turn}
\begin{itemize}
\raggedright
  \item \emph{Say it:} \emph{jônômôt}
  \item \emph{Say it:} \emph{jônômôt}, once more
  \item \emph{Say it:} \emph{jônômôt sôrkārer biruddhe.} — Public opinion is against the government.
  \item \emph{Recall:} say the Bengali for to protest, then the Bengali for a blockade, then say \emph{jônômôt} again
\end{itemize}

\begin{tcolorbox}[breakable,colback=teal!4,colframe=teal!35!black,title={Before you move on}]
What does \bn{জনমত} mean? (\textquotedblleft{}Public opinion\textquotedblright{}.) Say it once more.
\end{tcolorbox}
\section[\texorpdfstring{\bn{নীতি} (nīti)}{নীতি (nīti)}]{\bn{নীতি} (nīti) — a policy; a principle}

\label{lesson:BN-C349-niti}



\begin{quote}
Before the new one: say the Bengali for a blockade, then the Bengali for public opinion.
\end{quote}

\subsection*{You'll want to know: \bn{নীতি}}
\textbf{\bn{নীতি}} — \emph{nīti} — \textquotedblleft{}a policy; a principle\textquotedblright{}.

\begin{quote}
\bn{এই} \bn{নীতি} \bn{নিয়ে} \bn{অনেক} \bn{বিতর্ক} \bn{হল।}
\end{quote}

\begin{quote}
\emph{ei nīti niye ônek bitôrkô hôlô.}
\end{quote}

\begin{quote}
There was a lot of debate about this policy.
\end{quote}

\begin{cousinweb}
\textbf{\bn{নীতি}} (\emph{nīti}) is Sanskrit \emph{nīti}, guidance, wise conduct and policy (Monier-Williams). It is both a government's policy and a person's principle.
\end{cousinweb}

\begin{grammarlens}[title={What you've built}]
One more word for giving reasons and opinions.
\end{grammarlens}

\subsection*{Your turn}
\begin{itemize}
\raggedright
  \item \emph{Say it:} \emph{nīti}
  \item \emph{Say it:} \emph{nīti}, once more
  \item \emph{Say it:} \emph{ei nīti niye ônek bitôrkô hôlô.} — There was a lot of debate about this policy.
  \item \emph{Recall:} say the Bengali for a blockade, then the Bengali for public opinion, then say \emph{nīti} again
\end{itemize}

\begin{tcolorbox}[breakable,colback=teal!4,colframe=teal!35!black,title={Before you move on}]
What does \bn{নীতি} mean? (\textquotedblleft{}A policy; a principle\textquotedblright{}.) Say it once more.
\end{tcolorbox}
\section[\texorpdfstring{\bn{জনসংখ্যা} (jônôsôngkhyā)}{জনসংখ্যা (jônôsôngkhyā)}]{\bn{জনসংখ্যা} (jônôsôngkhyā) — population}

\label{lesson:BN-C349-jonosongkhya}



\begin{quote}
Before the new one: say the Bengali for public opinion, then the Bengali for a policy.
\end{quote}

\subsection*{You'll want to know: \bn{জনসংখ্যা}}
\textbf{\bn{জনসংখ্যা}} — \emph{jônôsôngkhyā} — \textquotedblleft{}population\textquotedblright{}.

\begin{quote}
\bn{আমাদের} \bn{গ্রামের} \bn{জনসংখ্যা} \bn{পাঁচ} \bn{হাজার।}
\end{quote}

\begin{quote}
\emph{āmāder grāmer jônôsôngkhyā pā̃ch hājār.}
\end{quote}

\begin{quote}
Our village has a population of five thousand.
\end{quote}

\begin{cousinweb}
\textbf{\bn{জনসংখ্যা}} (\emph{jônôsôngkhyā}) is a modern compound of Sanskrit \emph{jana}, people, and \emph{saṃkhyā}, a number or a count (Monier-Williams): the number of the people.
\end{cousinweb}

\begin{grammarlens}[title={What you've built}]
One more word for giving reasons and opinions.
\end{grammarlens}

\subsection*{Your turn}
\begin{itemize}
\raggedright
  \item \emph{Say it:} \emph{jônôsôngkhyā}
  \item \emph{Say it:} \emph{jônôsôngkhyā}, once more
  \item \emph{Say it:} \emph{āmāder grāmer jônôsôngkhyā pā̃ch hājār.} — Our village has a population of five thousand.
  \item \emph{Recall:} say the Bengali for public opinion, then the Bengali for a policy, then say \emph{jônôsôngkhyā} again
\end{itemize}

\begin{tcolorbox}[breakable,colback=teal!4,colframe=teal!35!black,title={Before you move on}]
What does \bn{জনসংখ্যা} mean? (\textquotedblleft{}Population\textquotedblright{}.) Say it once more.
\end{tcolorbox}
